\documentclass[11pt]{article}
\usepackage[hmargin=.9in,vmargin=.85in]{geometry}
\usepackage{amsmath,amssymb,amsthm,hyperref}
\hypersetup{colorlinks=true,linkcolor=blue,urlcolor=blue}
\newtheorem{theorem}{Theorem}
\newtheorem{lemma}[theorem]{Lemma}
\title{A finite bound for the number of unlocated balanced-array roots}
\author{ProgramComputer}
\date{November 26, 2026}
\begin{document}
\maketitle
\begin{abstract}
At height at least one million and between three and four rows per column, we isolate all but at most $\lceil n\sqrt{\log(100n)/(2d)}\rceil+4$ roots as simple negative real roots, for every row cap. The proof gives exact disk counts on circles centered at the origin. On each central circle a rotated seven-term approximation has uniformly positive real part. A monomial supplies the comparison near the two ends. In a million-row example, at least 501,926 of 508,000 roots are simple and negative.
\end{abstract}
Write $a_j$ for the number of labeled $n\times d$ binary matrices with all column sums $j$ and every row sum at most $r$, where $1\le r\le d$. Let $P(q)=\sum a_jq^j$. Its degree is $D=\lfloor nr/d\rfloor$ and each quota through $D$ is feasible: allocate the $d$ columns in consecutive blocks of $j$ cyclic row positions. Let $S(P)$ count its simple negative real zeros. Alexandersson asks whether all roots are real~\cite{original}.
\begin{theorem}\label{main}
If $n\ge10^6$ and $n/4\le d\le n/3$, then for every row cap,
\[
 D-S(P)\le \left\lceil n\sqrt{\frac{\log(100n)}{2d}}\right\rceil+4.
\]
\end{theorem}
This is a finite estimate with a numerical height and numerical constants. An earlier comparison with unrestricted root disks allowed a broader aspect range with unspecified constants~\cite{prior}. We obtain the count directly from the capped coefficients near one central term.

\section{A central disk}
Fix an internal quota $k$ with $x=k/n\in[1/8,7/8]$. Define
\[
 p_c=(k+1/2)/(n+1),\quad w_c=p_c/(1-p_c),\quad
 \eta=n\Pr\{\mathrm{Bin}(d,(k+3)/n)>r\}.
\]
\begin{theorem}\label{disk}
If $n,d$ satisfy Theorem~\ref{main} and $\eta\le1/100$, then $|q|<w_c^d$ contains exactly $k$ roots of $P$, counting multiplicity. The boundary contains no roots, and $(-1)^k P(-w_c^d)>0$.
\end{theorem}
\begin{proof}
Put $F_j=\binom nj^d$, $v=x(1-x)$ and $a=d/(2nv)\ge1/2$. For $|j|\le3$, write
\[
 T_j=F_{k+j} w_c^{dj}/F_k .
\]
We first show $|\log T_j+aj^2|<.001$. Let $w_0=x/(1-x)$ and $L(p)=\log(p/(1-p))$. For $j>0$,
\[
 \log\frac{\binom n{k+j}}{\binom nk}+j\log w_0
 =\sum_{i=1}^j\left[\log\left(1-\frac{i-1}{n(1-x)}\right)
 -\log\left(1+\frac{i}{nx}\right)\right].
\]
Replacing the logarithms by their linear terms costs at most $1216/n^2$ since $|\log(1+y)-y|\le y^2$ for $|y|\le1/2$ and $\sum_{i=1}^3((i-1)^2+i^2)=19$. The linear terms sum to $-j^2/(2nv)-j(1/2-x)/(nv)$. For negative $j$, reflect $k$ to $n-k$ and $w_c$ to $1/w_c$.

Now $p_c-x=(1/2-x)/(n+1)$, $L'(x)=1/v<10$ and $|L''(p)|\le162$ between $x$ and $p_c$ since both lie in $[1/9,8/9]$. Thus
\[
 \left|L(p_c)-L(x)-(1/2-x)/(nv)\right|\le86/n^2 .
\]
Multiplying the binomial and tilt errors by $d\le n/3$ gives $(1216+3\cdot86)/(3n)<.001$, proving the estimate.

For fixed quota $k+j$, choose columns independently as uniform $(k+j)$-subsets. Any one row has binomial load with parameters $d,(k+j)/n$. Monotonicity in this parameter and the union bound imply $1-\eta\le a_{k+j}/F_{k+j}\le1$. Consequently the seven capped terms, divided by the unrestricted central magnitude $F_kw_c^{dk}$, deviate in magnitude from $e^{-aj^2}$ by less than $.011e^{-aj^2}$. Here $.99e^{-.001}>.989$ and $e^{.001}<1.002$.

To make this comparison uniform on the circle, put $q=w_c^de^{i\phi}$ and rotate by $e^{-ik\phi}$. The real part of the ideal seven terms is
\[
 B_a(\phi)=1+2\sum_{j=1}^3 e^{-aj^2}\cos(j\phi).
\]
As a cubic in $u=\cos\phi$, its derivative is at least $2e^{-a}(1-4e^{-3a}-3e^{-8a})>0$, since $e^{3/2}>40/9$ and $e^4>30$. Its minimum is $1-2e^{-a}+2e^{-4a}-2e^{-9a}$, increasing in $a\ge1/2$ because $e^{3a}>4$. At $a=1/2$ the inequalities $e^{-1/2}<.607$, $e^{-2}>.135$ and $e^{-9/2}<.012$ give a strict lower bound $.032$. The sum of the ideal unsigned terms is below $2.51$ using $e^{-2}<.136$.

The unrestricted tilted sequence is log-concave. On either side the magnitude at distance three is less than $.013$ and its ratio to the term at distance two is below $e^{-5a+.002}<.09$. All subsequent outward ratios are no larger. The capped off-window terms are bounded individually by their unrestricted counterparts, so both tails together cost less than $2\cdot .013\cdot .09/(1-.09)<.003$. On the whole circle the real part of the normalized, rotated polynomial is therefore greater than
\[
 .032-2.51\cdot.011-.003>.001 .
\]
Its values stay in the open right half-plane. The original polynomial has winding number $k$ and the indicated sign on the negative ray. The argument principle gives the exact disk count.
\end{proof}

\section{Joining the circles}
At every internal index the orientation bound~\cite{orientation} gives
\[
 G_k=\frac{a_k^2}{a_{k-1}a_{k+1}}\ge
 \left[(1+1/k)(1+1/(n-k))\right]^d .
\]
When $\min(k,n-k)\le n/8$, it yields
\[
 \log G_k\ge\frac{d}{k+1}+\frac{d}{n-k+1}
 \ge \frac{n/4}{n/8+1}+\frac{n/4}{n+1}>2.2>\log9 .
\]
At the radius $u_k=\sqrt{a_{k-1}/a_{k+1}}$ the two immediate neighbors of the central term have relative magnitude $\tau=G_k^{-1/2}<1/3$. Log-concavity bounds all successive outward ratios by $\tau$. The total noncentral magnitude is below $2\tau/(1-\tau)<1$ times the central magnitude. The circle has no boundary roots, has count $k$, and has sign $(-1)^k$ on the negative ray.

Put $g=\sqrt{\log(100n)/(2d)}$. Whenever $k+3\le n(r/d-g)$, Hoeffding's binomial inequality gives $\eta\le n e^{-2dg^2}=1/100$. Use these central tests and the edge tests just proved. The uncovered internal indices belong to one consecutive block, with $k>nr/d-ng-3$. Because $D\le nr/d$, the block contains at most $\lceil ng\rceil+3$ indices. We use just the guaranteed tests for this assertion. If additional tests happen to pass they can only help.

Each certified disk has exactly its index as root count, so its radius is strictly ordered by index even where the formula for the radius changes. Supply index zero at the origin and index $D$ at a sufficiently large radius, using $P(0)=1$ and the leading sign. Each pair of consecutive certified indices brackets one negative real root, and their disk-count difference is one. It is consequently simple. If a block of $b$ indices was omitted, its enclosing annulus contains $b+1$ roots counted with multiplicity. This proves Theorem~\ref{main}, with the empty-block case giving all $D$ roots.

\section{A million-row example}
Take $n=10^6$, $d=250000$ and $r=127000$. Then $D=508000$ and $ng<6070$, so at least 501,926 roots are simple and negative. For $k=500000$, the binomial exponent $2d(r/d-(k+3)/n)^2$ exceeds $30$ and $e^{30}>10^{12}$. Hence Theorem~\ref{disk} gives exactly 500,000 roots strictly inside the unit disk. All these roots are simple and negative: the lower edge circles and the central circles up through this $k$ supply consecutive indices.

At radius $2^{-250000}$ the same example has exactly 333,333 roots strictly inside, again with no boundary root. Here $w=1/2$ and $k=333333$. The accompanying integer interval checker evaluates the seven unrestricted ratios $\bigl(\binom n{k+j}w^j/\binom nk\bigr)^d$ by outward-rounded repeated squaring. For $C_j=T_j+T_{-j}$ it certifies $C_1>4C_2+3C_3$ and hence that the real part of the seven unrestricted terms is minimized at $\phi=\pi$. The checker bounds the normalized minimum there minus both geometric tails and the cap loss by more than $.01$. The proof of Theorem~\ref{disk} then applies at this exact radius. It is smaller than the unit radius, so these roots are simple and negative as well. This count sharpens the radial interval obtained from a general finite estimate~\cite{distributions}.

\begin{thebibliography}{9}
\bibitem{original} P. Alexandersson. \href{https://www.symmetricfunctions.com/realRootedOpenProblems.htm\#conjBalancedArrays}{Balanced arrays of uniform matroids}. The symmetric functions catalog, accessed November 26, 2026.
\bibitem{prior} ProgramComputer. \href{https://github.com/ProgramComputer/open-math-research/tree/2cadc6002bda2bec41c982f5899cad820033074e/topics/balanced-arrays/dense-regime/below-cap}{Simple balanced-array roots below the cap}. Open Problems in Mathematics, 2026.
\bibitem{orientation} ProgramComputer. \href{https://github.com/ProgramComputer/open-math-research/blob/2cadc6002bda2bec41c982f5899cad820033074e/topics/balanced-arrays/orientation-normalization.md}{An orientation proof of normalized log-concavity for balanced arrays}. Open Problems in Mathematics, 2026.
\bibitem{distributions} ProgramComputer. \href{https://github.com/ProgramComputer/open-math-research/tree/f379722253bc3668d004a77df70fdf631de21273/topics/balanced-arrays/root-distributions}{Root distributions of balanced binary arrays with a common row cap}. Open Problems in Mathematics, 2026.
\end{thebibliography}
\end{document}
